% TOP_PROBLEM: {"id": 1432, "title": "Word-Representable Graphs: Forbidden Subgraph Characterization", "category_id": 3, "category": {"id": 3, "name": "graph_theory", "display_name": "Graph Theory", "description": "Problems involving graphs, networks, and their properties.", "slug": "graph-theory", "order_index": 3, "created_at": "2026-07-31T15:26:25.671Z"}, "status": "open"}
% FIRST_POSTED: null
% ENTRY_KIND: statement_only
% Imported from: batch07.tex; UnsolvedMath ID 1432
\documentclass[11pt]{article}
\usepackage[margin=25mm]{geometry}
\usepackage{fontspec}
\setmainfont{Latin Modern Roman}
\usepackage{amsmath,amssymb,mathtools}
\usepackage{unicode-math}
\setmathfont{Latin Modern Math}
\usepackage{xurl,hyperref}
\hypersetup{colorlinks=true,urlcolor=blue}
\setcounter{secnumdepth}{0}
\setlength{\parindent}{0pt}
\setlength{\parskip}{5pt}
\setlength{\emergencystretch}{3em}
\newcommand{\sref}[2]{\hyperlink{source:#1:#2}{[S#2]}}
\newcommand{\cataloguescope}[1]{\paragraph{Recorded target.}#1}
\begin{document}
% UnsolvedMath ID 1432; source catalogue code GRAPH-045
\setcounter{section}{1431}
\section{Word-Representable Graphs: Forbidden Subgraph Characterization}\label{1432}
{\small\sffamily UnsolvedMath ID 1432 · Graph Theory}
\subsection{Definitions and mathematical statement}

\paragraph{Shared notation.}$\mathbb N=\{1,2,\ldots\}$, and $\mathbb Z,\mathbb Q,\mathbb R,\mathbb C$ denote the integers, rationals, reals and complex numbers. Any other conventions are specified locally.


Let $G=(V,E)$ be a finite simple undirected graph. A word over $V$ is a finite sequence of vertices in which every vertex occurs. For distinct $x,y\in V$, let $w|_{\{x,y\}}$ be the word obtained by deleting every letter other than $x,y$. They alternate when this restriction has no equal consecutive letters, that is, it has the form $xyxy\cdots$ or $yxyx\cdots$, of either even or odd length. The word $w$ represents $G$ when, for every distinct $x,y$,
\[
 xy\in E\quad\Longleftrightarrow\quad x,y\text{ alternate in }w.
\]
Write $\mathrm{WR}(G)$ when such a word exists.

For $X\subseteq V(G)$, write $G[X]$ for the induced subgraph on $X$. A graph $F$ is a minimal forbidden induced graph for word-representability when $\mathrm{WR}(F)$ fails but $\mathrm{WR}(F[X])$ holds for every proper subset $X\subsetneq V(F)$. Word-representability is hereditary under taking induced subgraphs.
\paragraph{Statement recorded in the supplied source.}
Determine explicitly, up to isomorphism, the complete set $\mathcal F$ of minimal forbidden induced graphs for word-representability. A complete answer yields, for every finite simple graph $G$,
\[
 \mathrm{WR}(G)\quad\Longleftrightarrow\quad
 \text{no induced subgraph of }G\text{ is isomorphic to an }F\in\mathcal F.
\]
No finiteness assumption is imposed on $\mathcal F$. Merely defining $\mathcal F$ by minimal non-representability does not supply the requested explicit classification. \sref{1432}{2}
\subsection{Short English statement}
Explicitly identify every minimal graph that prevents word representation when it occurs as an induced subgraph.
\subsection{Sources}
\begin{description}
\item[\hypertarget{source:1432:1}{S1}] ULAM.AI and the UnsolvedMath Contributors, UnsolvedMath: A Curated Collection of Open Mathematics Problems, record 1432 (GRAPH-045). CC BY 4.0. \url{https://huggingface.co/datasets/ulamai/UnsolvedMath}.
\item[\hypertarget{source:1432:2}{S2}] Benny George Kenkireth, Gopalan Sajith and Sreyas Sasidharan, Word-Representability and Comparability: Minimal Forbidden Induced Subgraphs and Cover Number Bounds, arXiv:2502.06979v4 (2026), abstract and Section 1, pages 1--2. \url{https://arxiv.org/pdf/2502.06979v4}.
\end{description}
\label{end:1432}
\end{document}
